%\begin{acknowledgments}
\chapter*{Acknowledgments}
Un sentito e caloroso ringraziamento per questo lavoro va ai miei relatori, il prof. Stefano Giagu, il dott. Valerio Ippolito e il dott. Giuliano Gustavino.
Mi avete guidato, meeting dopo meeting, nella realizzazione di questo progetto, insegnandomi come cercare la nuova fisica e, soprattutto, l'approccio giusto per farlo, con spirito critico e curiosità.
Le lunghe chiacchierate in saletta, tra un caffè e un pranzo riscaldato, a discutere di tutti i dettagli, che oggi sembrano piccoli o persino ovvi, ma che si sono rivelati essenziali per la mia formazione.
Non so ancora cosa mi riserverà il futuro, ma so che, se sarò un buon fisico e ricercatore, sarà anche grazie a voi.\\
In secondo luogo, vorrei ringraziare Luca, per tutti i momenti di confronto e di supporto. 
Spero di tornare ad indossare presto la maglia rossa con il numero 16.\\
Ringrazio anche tutti i ricercatori e le ricercatrici \textit{senior} del gruppo ATLAS di Roma 1, per avermi accolto e condiviso con me le vostre esperienze e conoscenze.\\
Concludo i ringraziamenti ATLAS con un pensiero speciale ai membri giovani della saletta, laureandi e dottorandi, per aver condiviso con me le lunghe ed interminabili giornate in dipartimento, rendendole più leggere, piacevoli e divertenti, tra dubbi e curiosità, seri e meno seri.\\

%%%%%%%%%%%% OLD %%%%%%%%%%%%%%
% Un sentito ringraziamento va ai miei relatori, il prof. Stefano Giagu e la dott.ssa Graziella Russo, per avermi guidato nella realizzazione di questo lavoro. Vorrei inoltre ringraziarli per avermi accolto in \textit{saletta} e fatto conoscere questa bella e stimolante realtà.\\
% Giunti alla conclusione di questa dissertazione, applicando l'operatore $\hat{T}$, vorrei ringraziare:\\
% l'Associazione Italiana Studenti di Fisica (e non Fegato), AISF, per avermi dimostrato che c'è anche un'altra via per vivere la Fisica, un po' meno canonica ma ugualmente affascinante;\\
% il Pie$\check{z}$o per avermi mostrato l'aspetto umano dell'università e per avermi sopportato nei terribili periodi di sessione; \\
% Alessandro e Francesco, perché sono sufficienti tre birre e il solito bar per parlare di tutto e rendere i momenti difficili un brutto ricordo;\\ 
% il gruppo di Sanda Mareje, per essere stati parte viva della mia adolescenza e, oggi che siamo all'università, ogni volta che torniamo a Vasto, basta un messaggio per ritrovarci e sentirmi nuovamente a casa;\\
% Ale, che conosco da quando si frequentava lezione con un grembiule blu, con il quale ho condiviso tutti i principali eventi fino ad oggi, belli o meno, grazie della pazienza;
% la mia famiglia, per esserci stati nei momenti di difficoltà ed aver condiviso insieme momenti di festa.
% Infine, vorrei ringraziare i miei genitori, perché questo traguardo è anche vostro. Non deve essere facile vedere l'unico figlio partire di casa a 19 anni, direzione la capitale, con tanti sogni e paure. Non deve essere facile fare scelte e sacrifici per far vivere un giovane fuorisede. Non è stato facile, ma non mi avete fatto mancare mai il vostro supporto, fisicamente con la presenza di persona ed emotivamente, spronandomi a raggiungere vette sempre più alte. Grazie.\\
% Per concludere, vorrei riportare le parole che una persona cara mi aveva riferito in un'intervista. Alla mia domanda "Quale messaggio daresti ai giovani di oggi?" lui, rivolto a noi giovani, disse: "Ritrovate il colloquio e il dialogo con le persone. Ricominciate a guardare negli occhi chi vi circonda. Vi accorgerete che potrebbero esserci persone che condividono le vostre stesse passioni e voglia di costruire qualcosa di positivo. Fate le cose per passione e non tanto per farle."
%\end{acknowledgments}